\documentclass[12pt,a4paper,oneside]{article}
\usepackage[utf8]{inputenc}
\usepackage{times}
\usepackage[none]{hyphenat}
\usepackage{setspace}
\usepackage[protrusion=false,letterspace=75]{microtype}
\usepackage{ragged2e}
\usepackage{amsmath}
\usepackage{amsfonts}
\usepackage{amssymb}
\usepackage{makeidx}
\usepackage{graphicx}
\usepackage[left=2.5cm,right=2.5cm,top=2.5cm,bottom=2.5cm]{geometry}
\usepackage{color}
\usepackage[usenames,dvipsnames,svgnames]{xcolor}
\usepackage{changepage}
\usepackage{threeparttablex}
\usepackage{longtable}

\usepackage{cite}
	\renewcommand\citepunct{,}
\usepackage{url}
\usepackage[bookmarks=true,hyperindex,breaklinks]{hyperref}
	\hypersetup{bookmarksdepth=paragraph,colorlinks=true,linkcolor=Navy,urlcolor=DarkBlue,citecolor=DarkBlue}

%%%%%%%%%%%%%%%%%%%%%%%%%
\usepackage[authoryear,comma,round,semicolon]{natbib}
	\newcommand*{\doi}[1]{\href{http://dx.doi.org/#1}{doi: #1}}
	\setlength{\bibsep}{5pt} % plus 0.3ex
	%\setlength\parindent{10pt}
	\setlength\parindent{0pt}
	\setlength\bibhang{0em}
%\bibliographystyle{agsm}
%\bibliographystyle{plainnat}
%\bibliographystyle{abbrvnat}
%\bibliographystyle{unsrtnat}
\bibliographystyle{dcu}
%\bibliographystyle{apalike}
\usepackage{url}

\usepackage{soul}
%%%%%%%%%%%%%%%%%%%%%%%%%

\usepackage{float}
\usepackage[labelfont=bf,labelsep=period,font=singlespacing,skip=-10pt,hypcapspace=0pt,margin=5mm]{caption}
	\captionsetup[boxed]{skip=0pt,margin=5mm}
	\setlength{\belowcaptionskip}{-10pt}
\usepackage{graphicx}
	
\usepackage{listings}   
\usepackage{setspace}

\definecolor{codegreen}{rgb}{0,0.6,0}
\definecolor{codegray}{rgb}{0.5,0.5,0.5}
\definecolor{codepurple}{rgb}{0.58,0,0.82}
%\definecolor{backcolour}{rgb}{0.95,0.95,0.92}
\definecolor{backcolour}{RGB}{255,250,240}

\lstdefinestyle{mystyle}{
    backgroundcolor=\color{backcolour},  
    commentstyle=\color{codegreen},
    keywordstyle=\color{magenta},
    numberstyle=\tiny\color{codegray},
    stringstyle=\color{codepurple},
    basicstyle=\ttfamily\tiny,
    breakatwhitespace=false,         
    breaklines=true,                 
    captionpos=b,                    
    keepspaces=true,                 
    numbers=left,                    
    numbersep=5pt,                  
    showspaces=false,                
    showstringspaces=false,
    showtabs=false,                  
    tabsize=2
}

\lstset{style=mystyle}

\usepackage{array}
\newcolumntype{C}[1]{>{\centering\let\newline\\\arraybackslash\hspace{2pt}}m{#1}}

\usepackage{tabularx}
	\newcolumntype{Y}{>{\centering\arraybackslash}X}
\usepackage{caption} 
\captionsetup[table]{skip=10pt}
\usepackage{booktabs}
\usepackage{hyperref}
\graphicspath{{figures/}} % path to folder with figures
\usepackage{subcaption}
\usepackage{listings}
\usepackage{xcolor}
\usepackage{gensymb} % for \textdegree
\usepackage{array}
\usepackage{makecell}
\usepackage{multirow}
\usepackage{graphics}
\usepackage{color}
\definecolor{deepblue}{rgb}{0,0,0.5}
\definecolor{deepred}{rgb}{0.6,0,0}
\definecolor{deepmagenta}{RGB}{195,12,214}
\usepackage{multicol}
% ===============================================================================================
% start of the EDITOR section
% ===============================================================================================

% definitions

\def \newsection{\vspace{12pt}\textbf}
\def \newsubsection{\vspace{12pt}\textbf}
\def \newsubsubsection{\vspace{12pt}\textit}
\def \newpar{\vspace{6pt}}
\def \newref{\vspace{6pt}}
\def \figtab{\small}
\newcommand\tab[1][1cm]{\hspace*{#1}}

% first page header and dates information

\newcommand{\YEAR}{0}
\newcommand{\VOLUME}{0}
\newcommand{\NUMBER}{0}
\newcommand{\DOI}{10.2478/arsa-0000-0000}
%\setcounter{page}{1} 

% header

\usepackage{fancyhdr}
\lhead
{
	\begin{picture}(0,0)
	\put(0,-18){\includegraphics[width=2.7cm]{sciendo.png}}  
	\put(100,-6){\textbf{ARTIFICIAL SATELLITES, Vol. \VOLUME, No. \NUMBER ~-- \YEAR}}
	\put(100,-22){\textbf{DOI: \DOI}}
	\end{picture}
}

% footer

\lfoot
{
	\begin{picture}(0,0)
	\put(0,0){\includegraphics[width=2cm]{cr.png}}  
	\put(70,16){\tiny $\copyright$ 
	{\tiny The Author(s). This is an open-access article distributed under the terms of the Creative Commons Attribution License (CC BY 4.0,}}
	\put(70,8){\tiny {\href{https://creativecommons.org/licenses/by/4.0/}{https://creativecommons.org/licenses/by/4.0/}), which permits use, distribution, and reproduction in any medium, provided that the}}
	\put(70,0){\tiny Article is properly cited.}
	\end{picture}}
\cfoot{}
\renewcommand{\headrulewidth}{0pt}
\renewcommand{\footrulewidth}{0pt}

% -----------------------------------------------------------------------------------------------
% end of the EDITOR section
% -----------------------------------------------------------------------------------------------

\graphicspath{{figures/}} % path to folder with figures

\begin{document}
\sloppy
\thispagestyle{fancy}

\vspace*{38pt}

% ===============================================================================================
% start of the AUTHOR section
% ===============================================================================================

% title A script-driven approach to mapping satellite-derived topography and gravity data over Zagros Fold and Thrust Belt, Iran

\begin{center}{\large\bf\uppercase{
Machine Learning Methods of Remote Sensing Data Processing for Mapping Salt Pan Crust Dynamics in Sebkha de Ndrhamcha, Mauritania}}
\end{center}

\vspace*{4pt}

% author(s)

\begin{center}
\textls{
Polina \uppercase{Lemenkova}$^1$  \\

\vspace{6pt}

$^1$Alma Mater Studiorum -- Università di Bologna, Dipartimento di Scienze Biologiche, Geologiche ed Ambientali, Bologna, Emilia-Romagna, Italy

%$^2$Università di Bologna, Dipartimento di Scienze Biologiche, Geologiche e Ambientali, Bologna, Italia.

\vspace{6pt}

e-mails: polina.lemenkova2@unibo.it
}
\end{center}

\vspace{30pt}

% abstract and keywords

{\bf ABSTRACT.} The advances in Machine Learning (ML) and computer technologies enabled to process satellite images using programming. Environmental applications that handle Remote Sensing (RS) data for spatial analysis use such an approach, for example, Python’s library scikit-learn using algorithms on pattern identification, predictions or image classification. This paper presents an ML method of satellite image processing using Geographic Resources Analysis Support System (GRASS) Geographic Information System (GIS). The aim is to classify multispectral Landsat images using ML for identification of changes in salt pans of West Mauritania, Africa over the period 2014–2023. We define 10 classes of land cover categories and perform analysis of geological, lithological and landscape setting, and then introduce the principles, algorithms and processing of the ML methods of GRASS GIS. The following classification models were employed to implement image classification with training: Random Forest (RF), Decision Tree, Gradient Boosting and Support Vector Machine (SVM). The results were compared with clustering performed by k-means and maximum likelihood discriminant analysis. The cartographic visualisation and validation was implemented through accuracy analysis. Results for the best performing SVM model with seven-band input produced an overall accuracy of 76\%, for the RF model - 73\% , compared to 69\% for Decision Tree Classifier -69\% and for Gradient Boosting Classifier - 67\% The SVM model embedded in GRASS GIS generates robust land cover maps with good accuracy from multispectral satellite images. The paper demonstrated an ML-based automated approach to satellite image processing, which links Artificial Intelligence (AI) with cartographic tasks.

\newpar {\bf Keywords:} Machine Learning, Random Forest, Support Vector Machine, GRASS GIS; Satellite Image, Remote Sensing, Saline Lake.

%\pagebreak
\newsection {1. INTRODUCTION}

\subsection*{1.1. Background}

\newpar During recent decades, the landscapes of the Earth have been changing at an unprecedented rate due to the cumulative effects from climate change and human activities. Mauritania is one of the African countries most affected by the cumulative effects of natural and anthropogenic effects on land patterns  \citep{CHOPLIN201749}. Given the challenges from the environmental and climate change as well as water scarcity in arid areas of Africa, it is becoming more and more crucial to comprehend, map and monitor the sub-Saharan landscapes better for effective land management. Predictive ecosystem modelling and appropriate land management policies in arid regions of West Africa strongly depend on computational data, which support environmental modelling in the context of land cover change \citep{BOJER2024121394,PILLAY2021106595,Lemenkova202414}. In environmental mapping of Africa, digital mapping to analyse landscape dynamics remains an open problem despite selected existing research \citep{ALI2024113123,WALTHER2021146093,TWUMASIANKRAH2024176219,Lemenkova202413}. 

\newpar Mosaics of vegetation, cropland and shrubland interspersed with sandy areas have historically played significant role in the landscape mosaic of West Africa. However, from the middle of the 20th century, land cover structure in Mauritania has undergone substantial alterations, which were brought about by several triggers. The main drivers of ecosystem dynamics in Africa remain natural factors, climate change and human activities. Natural factors that cause land cover changes in Africa include a variety of processes. Some of them are geological hazards, geomorphological instability in sandy and rocky regions that cause soil degradation, climate-related issues along the coasts of Atlantics (floods, droughts, fluctuations in precipitation and temperature that affect coastal vegetation) and diversity of soil types. To understand the complexity of these factors, advanced computational modelling methods are required.

\newpar Within the scope of ecosystem research, there has been some progress towards computational modelling that mostly uses traditional methods of Geographic Information System (GIS) \citep{DUNGA2024108989,EBRAHIM2024101374,SHOKO2024103796,Lemenkova202412,Holloway}. This paper presents a Machine Learning (ML) framework to direct the analysis of satellite images for detecting and mapping changes in land cover in western Mauritania with a case of saline lake. Comparing several ML methods incorporated into the workflow of Remote Sensing (RS) data processing and cartographic scripts, the changes in fluctuating saline lake were detected using classified satellite scenes.

\newpar By utilising Earth observation (EO) data, land cover change detection aims to visualise and monitor the dynamics of various patterns of the landscape as a linked human–environment system. For advanced monitoring, one of the important approaches for identification landscape dynamics in land cover types is the analysis of changes and the rate of exposure of landscapes to climate and environmental effects using multitemporal geospatial RS data. ML is widely recognised as advanced computer application for processing RS data. Besides, it has a powerful functionality as a branch of Artificial Intelligence (AI) \citep{YANG2024102087,RATHORE2016122,LEMENKOVA2025100180,CEDRIC2022100049}, including application of RS data processing. The advantages of the RS data for environmental and geological mapping have been discussed in many previous works (e.g. \citep{Parajuli,Niang2022,Ashpole}). Among the benefits of the RS data, such as satellite images \citep{GRANDJEAN2006220} , is mentioned their cost-effectiveness for geological exploration, since the RS data enable precise coverage of vast regions, such as Africa regions, with essential data for recourse exploration. The need for satellite images for natural resources management gives rise to the question of their effective processing.

\newpar With the recent advances of computing and programming technologies, more and more RS data, namely satellite images, can be processed using ML \citep{BELMAHDI2023100906,Lemenkova202021,GOMEZ201929}. Correspondingly, new types of methods of satellite image processing based on ML emerged recently. Land cover mapping from RS data has progressed from visual and simple statistical approaches to advanced image recognition techniques. These include algorithms developed in Python scikit-learn library and include, for instance, Artificial Neural Networks (ANNs), Random Forest (RF), Convolutional Neural Networks (CNNs), Support Vector Machines (SVMs) and many more. Existing reports on the use of these methods prove their effectiveness in computational RS \citep{Lemenkova202411,GAWUSU2024100065,Lemenkova202410,NENKAM2024117007}. The key feature and advantage of ML consists in its capability to imitate intelligent human behaviour in the algorithms of image processing and pattern recognition \citep{EBRAHIMY2023103390,PANDE2024141035,Lemenkova202217}. The ML-based approaches are typically distinguished from conventional GIS methods of image classification and processing by their advanced data handling consisting of application of neural networks for geospatial analysis. 

\newpar ML methods can be applied for pattern detection and categorisation of land cover types, which could be regarded as a new technical tool of satellite image processing in cartography and geoinformatics. ML-based image processing, as in the case of environmental mapping, can be integrated into various applications of land management. For example, climate-related events such as droughts of floods may change existing landscape patterns and make the current maps of land cover types outdated. In particular, at the current stage, satellite images can be processed by ML methods for supervised classification using training data and machine logic, as pixel based \citep{RAHMAN2020100410,AYOADE2023337} or object based \citep{KARILA2023100046,Lemenkova202024,Lemenkova201528}. The advantages of both approaches were discussed in details in earlier review works \citep{NASIRI2023103555,ABUJAYYAB2023101802}. 

\newpar An overview of the field demonstrates existing case studies which use independent specialised packages for image processing and classification, as well as integrated approaches supporting GIS through programming. For instance, many programming libraries and packages of well-known languages (e.g. Python or R; \citep{Lemenkova202324} attempt to evaluate patterns on satellite images using geospatial libraries \citep{geomatics5010005,lemenkova2024approche}, and some of the advanced software consider ML approaches and scripting techniques. Among these, the open-source software Geographic Resources Analysis Support System (GRASS) GIS presents effective approach \citep{NETELER2012124,MitasovaNeteler}. The available GRASS GIS tools in the ML area (‘r.learn.train’, ‘r.learn.predict’) are not implemented natively, but rather depend on such methods which implement the algorithms used in this study. Straightforwardness and robustness of algorithms embedded in GRASS GIS are two fundamental advanced features that support the progress of script-based cartographic methods and allow its development towards automation, reliability and flexibility of mapping \citep{ROCCHINI201382,Lemenkova202307,ROCCHINI2017166}. Besides, recent version of this software (GRASS GIS 8.4) has been updated for functionality of satellite image processing. Extended functions now enable to reuse spectral signature files of classifiers, improved output of accuracy assessment and selected methods of ML classifiers \citep{Nartiss}. Finally, this software supports diverse approaches of ML through different algorithms combining supervised learning.

\subsection*{1.1. Objectives and Motivation}

\newpar In the dry areas of Mauritania, salt pan habitats have been identified as important suppliers of mineral elements. Nevertheless, salt pans are also sources of dust emissions, which have effects on the climate processes, lead to the degradation of groundwater, contribute to soil and ocean fertilisation, as well as have impacts on human and ecosystem health in coastal Mauritania. Dust type and emissivity, specifically for Sebkha de Ndrhamcha, are related to the mineralogical and physical crust composition of soil and local geological features. Understanding seasonal fluctuations of the salt pan and their effects on regional environmental conditions is, therefore, crucial for land management in the arid climate of Mauritania. The advances in modelling salt dynamics can be used, for instance, for predicting dust emissions, which depends on the specifics of the salt pan surface and its changes over time. Knowing how exactly salt pan changes over time may be used for prediction of water scarcity and protection of plants in agricultural activities of farmers.

\newpar Hence, there is a growing demand from the environmental decision makers for modelling salt pans in the arid regions of West Africa. Nevertheless, the vast area and challenging access in many regions of Mauritania limit the potential of typical field approaches to investigate salt pan surfaces. At the same time, accurate mapping in the areas that are hard to reach presents difficulties and requires machine-based methods and RS data. These issues have been addressed by processing of RS data because of their capacity to give multitemporal observations and to cover wide areas. To process the RS data effectively, the cartographic use of ML tools presents a novel approach for automated mapping of salt pans. This is principally a new step in mapping arid regions of West Africa, which presents the main goal of this study: application of the advanced ML tools for mapping seasonal fluctuations of the salt pans in Mauritania.

\newpar Nevertheless, due to the specific spatiotemporal setting of satellite images, GIS methods are different from scripting libraries related to general programming, and there is a lack of efficient reports on applications of ML to satellite image processing. Thus, there is still a lack of efficient reports on mapping salt pans in the western region of Mauritania which evaluate changes using satellite images according to the review on related works on West Africa and coastal Mauritania. In view of this, we apply and propose a novel method of satellite image classification using GRASS GIS, using the combination of ‘r.learn.train’, ‘r.learn.predict’ and ‘r.category’ modules applied for RS data processing with a case of Landsat images. Instead of a single method, we tested a group of algorithms and evaluated their performance for cartographic data processing and image classification. Various algorithms of ML can be used for environmental monitoring by combining time series analysis of several satellite scenes within the GRASS GIS tools. Furthermore, it is easy to implement and maintain the scripts using the presented codes commented in the ‘Materials and methods’ section. In this way, this manuscript contributes to the future reuse and applications of the proposed methods in similar studies on satellite image processing for environmental monitoring.

\newpar This study aims at using ML techniques to study sebkha of Mauritania to test the usefulness of several ML algorithms for processing multispectral Landsat satellite data. To this end, data in the short time period (2014–2023) were applied for analysing pan surface properties of Sebkha de Ndrhamcha and its changes over time. Since no similar studies that use ML tools and RS data for mapping Sebkha de Ndrhamcha exist in modern literature, this study presents an interdisciplinary approach that uses advanced technical tools for environmental purposes

\newpar More specifically, in this paper, we present the use of the advanced ML modules for image classification using GRASS GIS  \citep{GRASS_GIS_software} with the case of Mauritania, West Africa. Specifically, we employ, compare and relate several functions of ML to image processing, namely RF Classifier, Decision Tree Classifier (DTC), Gradient Boosting Classifier (GBC), SVM Classifier, evaluate their performance and compare the results with the approaches of the conventional methods of image classification, such as Minimal–Maximal Discriminant Analysis and k-means clustering. The potential of the ML modules of GRASS GIS \citep{NetelerMitasova2002} for mapping ecological patterns and processes in arid climate of Sahel is exemplified with the case of fluctuating sebkhas of coastal Mauritania and associated landscape dynamics in the areas of salt dunes of sub-Saharan Africa.

\subsection*{1.3. Structure}

\newpar The remainder of the paper is organised as follows. The section ‘Study Area’ gives a formal description of the study area and related geological, topographic and environmental setting illustrated by thematic maps. The section ‘Materials and Methods’ describes the data used for image processing with mentioned metadata, proposes the approaches of GRASS GIS including the comments on its ML module, their principles of RS data processing, functionalities, structure of methodology, and comments the scripts in Appendix. ‘Results and Discussion’ discusses several outcomes of image processing by examples of ML-based maps. The section ‘Conclusion’ shows the implementation of the ML methods of GRASS GIS applied for the salt lakes of Mauritania and concludes the manuscript.

\newsection {2. STUDY AREA}

\newpar The study area is located in the coastal region of West Mauritania, which includes the salt pan of the Sebkha de Ndrhamcha \ref{fig01}. \newpar Physical processes affecting the formation of sebkhas include the sedimentation in coastal, fluvial and lacustrine environments, as well as aeolian-dominated conditions (the strength and direction of winds). For instance, the Atlantic continental margin to the west of Sahara presents one of the largest dust sources on Earth and is strongly influenced by mass wasting processes \citep{HENRICH2010178}. Huge aeolian dune fields affect the sedimentary regime on the continental slope and surrounding sebkhas.

\newpar The prevalence of one or more of these settings favours formation of the salt pans and results in different level of their salinity \citep{TARI2017331}. Thus, marine coastal sabkhas are more associated with the processes of evaporation and accumulation of marine deposits, while continental playas or inter-dune sabkhas are related to barchan formation and dune mobility \citep{HANDFORD1981255,Ould}. Different salinity of salt lake sediments enables to distinguish the north-west Africa hypersaline chotts with saline levels surpassing that of ocean water due to high concentrations of brines \citep{Lemenkova202313} and non-vegetated sebkhas with moderate salinity \citep{Lemenkova202322} from vegetated sebkhas with a relatively low salinity that ranges from 2 to 42 g NaCl per litre \citep{SOULIEMARSCHE200869}.

\newpar The geological setting of western Mauritania is presented by a series of Quaternary marine deposits formed in the gulfs of Senegalo–Mauritanian Tertiary–Cretaceous basin (Figure \ref{fig02}). During the uplift of the basin in Miocene, this area had undergone four successive transgressions and three regressions, which resulted in the accumulated continental formations covering marine deposits with sandy-glauconitic facies and diatomite layers \citep{barry2003contribution}. Several periods of eustatic movements of the world’s oceans formed the current Quaternary shape of the coastal zone in West Mauritania, with lakes and lagoons becoming sebkhas and bays becoming lagoons during the alternating periods \citep{EINSELE1974282,Elouard}.

\begin{figure}[H]
\centering
	\includegraphics[width=16.0 cm]{Fig_01.jpg}
	\vspace*{20pt}\caption{Topographic map of Mauritania showing the extent of the study area of the Sebkha de Ndrhamcha salt pan (rotated yellow square). Data: General Bathymetric Chart of the Oceans (GEBCO)/Shuttle Radar Topography Mission (SRTM). Software: Generic Mapping Tools (GMT). Map source: author.
	\label{fig01}}
\end{figure}

\begin{figure}[H]
\centering
	\includegraphics[width=16.0 cm]{Fig_02.jpg}
	\vspace*{20pt}\caption{Surficial geology and lithologic units of Mauritania. Data: USGS, GEBCO. Software: QGIS. Map source: author.
	\label{fig02}}
\end{figure}

\newpar Successive positive epeirogenic episodes resulted in reduced size of the marine gulfs off Nouakchott since the early Pleistocene to Holocene. The Tafaritian sandstones of the Aaiun-Tarfaya Basin belong to the continental environments, such as the lacustrine and sebkha deposits formed under shallow or aerial conditions of Saharan–Sahel environments. Subsequently, the former large marine Tafaritian embayment with uplifted sediments was replaced by the sparse lakes and sebkhas \citep{GIRESSE2000123,GIRESSE1988241}. Currently, the mineral content of coastal sabkhas consists of a combination of terrigenous clastics, carbonate–sulphate (or gypsum) minerals, soluble salts (e.g. halites) and similar types of sediments \citep{Kocurek}. Such specific mineral content enables discrimination of these sediments on the satellite images using spectral reflectance of the saline minerals which differ by colours and brightness from the surrounding areas and sands. The lacustrine deposits discovered in other regions of Sahara gave rise to the notion of ‘Saharian Seas’ and current lakes as remnants of the endoreic depressions reflecting past geological processes \citep{PLAZIAT1991383,Lemenkova202323}.

\newpar The origin of the Mauritanian sabkhas is strongly related to the effects of climatic and geological factors. The diversity of geological setting of Mauritania geological history is characterised by the great variety of rocks and geological structures. Among them, the Reguibat Shield (500 by 2,000 km, stretching SW to NE) forms the northern part of the West African Craton and belt of Proterozoic rocks (Archean, Paleoproterozoic to Neoproterozoic ages) (Figure \ref{fig03}). The plate tectonic movements formed the West African Shield, which borders the Taoudeni basin \citep{HAMOUD2021415}. During the Holocene, in the area located on West African Shield and Ougarta Uplift (Figure \ref{fig03}) prevailed wet climate with repetitive marine transgressions and climatic oscillations (succession of humid/arid phases in the region of modern Nouakchott) \citep{einsele1974holocene}. Such processes are now reflected in the Holocene stratigraphy and geomorphological landforms \citep{Wissmann}. This is also revealed in the palynological sequences according to stratigraphical and pollen analysis \citep{MEDUS1987167}. Besides, the fluctuations in sea water inundation resulted in continuous calcification of the evaporite deposits and marine sulphates that increased the salinity of the area \citep{MANNINGBERG2024111974}.

\begin{figure}[H]
\centering
	\includegraphics[width=14.0 cm]{Fig_03.jpg}
	\vspace*{20pt}\caption{Geological provinces around Mauritania. Data: USGS, GEBCO. Software: QGIS. Map source: author.
	\label{fig03}}
\end{figure}

\newpar The climate factors are expressed in the temperature, evaporation and drought in the region, while physical processes determine the morphology of sebkhas and the sedimentary response to those processes. Moreover, the effects of climate-oceanological processes on the coastal habitats can be illustrated through the seasonal upwelling, which affects the coastal productivity and distribution of nutrients for aquatic habitats \citep{VAZQUEZ2023166391}. Such cyclic circulation of water layers has a great ecological and socio-economic importance for the coastal region of Mauritania.

\newpar The distribution of minerals, geomorphic landforms and types of soil strongly affect the overall mosaic of habitats and land cover types; they also strongly affect the agriculture resources and crop management \citep{QUEBEDEAUX1984133}. For instance, in coastal areas, tidal flats and shallow waters favour the distribution of aquatic and regularly flooded vegetation \citep{WIJNSMA1999233}, while sandy desert regions are mostly covered by sparse herbaceous vegetation. The general map of land cover types and their distribution over Mauritania is presented in Figure \ref{fig04}.

\begin{figure}[H]
\centering
	\includegraphics[width=16.0 cm]{Fig_04.jpg}
	\vspace*{20pt}\caption{Land cover types in Mauritania. Data: FAO, OpenStreetMaps. Software: QGIS. Map source: author.
	\label{fig04}}
\end{figure}

\newpar Similar to other regions of Sahara with salt lakes (e.g. \citep{ALI20231}), the formation of salt pans in Mauritania is related to salt tectonics and accumulation of evaporite deposits. Besides strong climatic factors, the impact of geological setting on the formation of salt pans is explained by rift geometries that are defined by horsts and tilted fault blocks and generate base-salt flat topography, which affects salt accumulation, flow direction, diapirism and deformation of overburden layers \citep{Pichel}. The distribution of salt minerals is concentrated in desert areas with highly arid climate where evaporation exceeds precipitation. Therefore, the formation of salt lakes, salt pans, chotts and sebkhas in Africa is mostly located in northern regions close to Sahara \citep{VANASTEN200313,ALVARO2012209}. The crust of salt deposits and salty compacted sands creates natural resources for the country where salt is being produced for market and industrial purposes. 

\newpar Current environmental problems in the region include risks of flooding in Nouakchott, which are due to a combination of both natural factors (wave actions and wind strength) \citep{semega2008energie} and human activities (infrastructure built along the coasts and fragility of the barrier beach) \citep{Senhoury2016,ould2007risques}. For instance, the semi-arid climate of Mauritania creates additional constraints for agricultural activities, which results in construction of large-scale irrigation schemes aimed at increasing yield and irrigation intensity \citep{BORGIA20121}. Nevertheless, unsustainable crop management is a threat to ecological sustainability of coastal areas and contributes to the salinisation of soils. Lack of effective drainage system for rain and marine waters and a sewage sanitation system increases the problem of salt accumulation along the coasts of Mauritania and creates additional factors for soil salinisation around Nouakchott.

\newpar Flat topographic landforms of the Mauritanian coasts along the Atlantic Ocean, such as the extent of the Aftout Es Sahli (a geomorphological depression running parallel to the coastline) and Sebkha N’Dramcha (a salt pan continuing the Senegalese–Mauritanian sedimentary basin to the north off Nouakchott), create additional risks for the occasional inundations and floods \citep{Vermeer2010}. Sea water intrusion contributes to the increase of salinity in the coastal areas and supports formation of sebkhas and accumulation of salt in the submerged landscape patches. Another notable characteristic of the Mauritanian coasts is the upwelling tidal system, which is formed by the effects of the coastal slope topography, ocean currents and winds. As a result, the coastal region of Mauritania is characterised by almost permanent upwelling with maximal intensity in boreal winter and spring \citep{Versteegh}.

\newpar The cumulative effects of human activities and climate processes, briefly discussed above, lead to the decline of groundwater resources in Mauritania, which is especially notable for its semi-arid areas \citep{MOHAMED2017280}. All these factors lead to the crucial problems of water deficits and their associated shortages in Mauritania, such as the need for water desalination and the problem of water supply \citep{BAYODRUJULA201099}, and contribute to the formation of the salt pan Sebkha de Ndrhamcha.

\newsection {2. MATERIALS AND METHODS}

\newpar The data used in this study include six satellite multispectral Landsat images by National Aeronautics and Space Administration (NASA), whose main properties are summarised in Table \ref{tab1}. 

\begin{table}[H]
\scriptsize
\caption{Metadata of the multispectral satellite images Landsat 8-9 OLI/TIRS, used in this study, obtained from USGS \textsuperscript{1}.\label{tab1}}
\begin{tabularx}{\textwidth}{cccccc }
			\toprule
			\textbf{Date} & \textbf{\makecell{Spacecraft /\\ID}} & \textbf{\makecell{Path/\\row}} & \textbf{Entity product ID} & \textbf{Scene ID} & \textbf{\makecell{Cloud/\\Coverage}} \\
			\midrule
			27/04/2014 & Landsat 8 & 205/47 & \makecell{LC08\_L2SP\_205047\_20140427\_20200911\_02\_T1} & LC82050472014117LGN01 & 0.00 \\
			03/04/2017 & Landsat 8 & 205/47 & \makecell{LC08\_L2SP\_205047\_20170403\_20200904\_02\_T1} & LC82050472017093LGN00 & 0.04 \\
			22/04/2018 & Landsat 8 & 205/47 & \makecell{LC08\_L2SP\_205047\_20180422\_20201015\_02\_T1} & LC82050472018112LGN00 & 0.04 \\
			11/04/2020 & Landsat 8 & 205/47 & \makecell{LC08\_L2SP\_205047\_20200411\_20200822\_02\_T1} & LC82050472020102LGN00 & 0.17 \\
			09/04/2022 & Landsat 8 & 205/47 & \makecell{LC09\_L1TP\_205047\_20220409\_20230422\_02\_T1} & LC92050472022099LGN01 & 0.00 \\
			28/04/2023 & Landsat 9 & 205/47 & \makecell{LC09\_L2SP\_205047\_20230428\_20230430\_02\_T1} & LC92050472023118LGN00 & 0.01 \\
			\bottomrule
		\end{tabularx}
	\noindent{\footnotesize{\textsuperscript{1} The Sensor ID is common for all the scenes: Landsat Operational Land Imager and Thermal Infrared Sensor (OLI/TIRS), Collection 2 Level-2. Image courtesy of the US Geological Survey (USGS). Product DOI: \href{https://doi.org/10.5066/P9OGBGM6}{10.5066/P9OGBGM6}.}}
\end{table}

\newpar The input raster scenes were based on the six Landsat images. The imagery data of the Landsat 8-9 Operational Land Imager and Thermal Infrared Sensor (OLI/TIRS) satellites were used in this study because they offer systematic coverage of the study with minimal cloudiness. Besides, technical characteristics of the radiometric and geometric features of the scenes are suitable for environmental monitoring in short-term period (2014–2023) using pixel-based methods of image processing. Such tools are applicable in GRASS GIS which includes a wide range of automated algorithms of image classification in its libraries. 

\newpar Image analysis and classification of the Landsat 8-9 OLI/TIRS satellite images using pixel-based clustering method was done by k-means and an unsupervised method, which is included in the maximum likelihood discriminant analysis classifier by GRASS GIS. Pixel-based image classification techniques, such as k-means, are becoming more and more acknowledged as a useful cartography method for RS data processing to take advantage of the vast amounts of spaceborne data. Such methods enable creating representative maps for annual or seasonal time periods to use either as input for model building or as visualisation of the environmental situation in the given period. The raw images are presented in Figure \ref{fig05}, which shows the scenes for April month of years 2014, 2017, 2018, 2020, 2022 and 2023.

\begin{figure}[H]
	\begin{subfigure}[b]{.35\textwidth}
		\centering
			\includegraphics[width=0.87\textwidth]{Fig_05a.jpg}
		\caption{2014}
	\end{subfigure}%
	\begin{subfigure}[b]{.35\textwidth}
		\centering
		\includegraphics[width=0.87\textwidth]{Fig_05b.jpg}
		\caption{2017}
	\end{subfigure}%
	\begin{subfigure}[b]{.35\textwidth}
		\centering
		\includegraphics[width=0.87\textwidth]{Fig_05c.jpg}
		\caption{2018}
	\end{subfigure}%
\\
\vfill \vspace{1mm}
	\begin{subfigure}[b]{.35\textwidth}
		\centering
			\includegraphics[width=0.87\textwidth]{Fig_05d.jpg}
		\caption{2020}
	\end{subfigure}%
	\begin{subfigure}[b]{.35\textwidth}
		\centering
		\includegraphics[width=0.87\textwidth]{Fig_05e.jpg}
		\caption{2022}
	\end{subfigure}%
	\begin{subfigure}[b]{.35\textwidth}
		\centering
		\includegraphics[width=0.87\textwidth]{Fig_05f.jpg}
		\caption{2023}
	\end{subfigure}%
\vspace*{20pt}\caption{Landsat 8-9 satellite images covering salt pan Sebkha de Ndrhamcha in western Mauritania in natural colours showing floodplain for 6 years (always April): (\textbf{a}) 2014, (\textbf{b}) 2017, (\textbf{c}) 2018, (\textbf{d}) 2020, (\textbf{e}) 2022, (\textbf{f}) 2023. Brief explanations of the colours: black colour indicates water areas of the Atlantic ocean, light beige colour represents sandy areas, brown colour indicates desert and bare lands, dark brown colour shows wet areas, grey colour indicates artificial surfaces (roads and urban areas) and cyan colour indicates salt pans on the original images. The colours are in Red Green Blue (RGB). Data source: US Geological Survey (USGS), downloaded from the EarthExplorer repository. }\label{fig05}
\end{figure}

\newpar The map in Figure 1 shows the topographic setting and the location of the study area. The map is made using the Generic Mapping Tools Version 6.4.0 scripting toolset \citep{Wessel}. The programming scripts used for cartographic techniques are described in earlier works (e.g. \cite{Lemenkova202203,Lemenkova202216,Lemenkova202212}).

\newpar The images were processed using the scripts provided in Appendix. First, the data were imported into the GRASS GIS environment and preprocessed using the script provided in Listing \ref{lst01}. Then, the images were classified using unsupervised clustering methods with k-means algorithms. The GRASS GIS script used for unsupervised classification is provided in Listing \ref{lst02}. The results of this classification were used as training polygons for ML methods, providing the ‘seed’ information for supervised learning. The next step included the application of RF Classifier for image classification. The programming code for this is provided in Listing \ref{lst03}. The following step included the use of the DTC for image classification, with the script demonstrated in Listing \ref{lst04}. The algorithms of GBC for image classification were performed using the programming code shown in Listing \ref{lst05}. 

\newpar Finally, the algorithm of SVM Classifier was applied for supervised image classification using the script in Listing \ref{lst06}. The performance of the algorithms demonstrated different time consumption and machine resources. Therefore, to note the difference in speed and the time required for execution, we summarise the performance of the methods in Table \ref{tab2}. The quickest algorithm was the DTC, while the longest period of data processing was demonstrated by the SVM Classifier. The data were processed using the MacOS 14.2.1 (23C71), Apple M1 Chip Processor.

\begin{table}[H]
\small
\caption{Processing time of the satellite images Landsat-8 OLI/TIRS showing the effectiveness of ML methods executed by the GRASS GIS \textsuperscript{1}.\label{tab2}}
\begin{tabularx}{\textwidth}{cp{4cm}}
			\toprule
			\textbf{Method} & \textbf{Processing Time} \\
			\midrule
			Clustering & $<$ 1 minute \\
			Min-max discriminant analysis & ca. 25 sec \\
			Random Forest Classifier & 9 minutes  \\
			Decision Tree Classifier & ca. 34 sec \\
			Gradient Boosting Classifier & 23 minutes  \\
			Support Vector Machine Classifier & 47 minutes  \\
			\bottomrule
\end{tabularx}
	\noindent{\footnotesize{\textsuperscript{1} The Sensor ID is common for all the scenes: Landsat OLI/TIRS (Operational Land Imager and Thermal Infrared Sensor), Collection 2 Level-2. Image courtesy of the U.S. Geological Survey (USGS). Product DOI: \href{https://doi.org/10.5066/P9OGBGM6}{10.5066/P9OGBGM6}.}}
\end{table}

%\pagebreak
\newsection {3. RESULTS AND DISCUSSION}

\newpar Environmental and hydrological planning in arid regions of Africa is essential and is largely based on using geoinformation such as satellite images. Due to their severe climate and vast spatial regions, salt pans are extremely dynamic natural features of Mauritania that are challenging to investigate using in situ approaches. This study employs spaceborne multitemporal (2014–2023) multispectral Landsat data combined with thematic on topography and geology of Mauritania to document the spatial distribution of surface crust types over time in the western part of the country: region of Sebkha de Ndrhamcha. For this purpose, six satellite images were collected from the US Geological Survey (USGS) and analysed for spectral differences covering the period 2014–2023, which reflects different seasons and surface conditions of the salt pan. The data regarding the mineralogical and geological setting were analysed using the thematic maps (Figures \ref{fig02} and \ref{fig03}) as well as distribution of land cover types according to Food and Agriculture Organization (FAO) (Figure \ref{fig04}).

\newpar An approach based on the GRASS GIS tools was developed to perform the spatiotemporal analysis of the salt pan crust dynamics in a dense time series consisting of six Landsat 8-9 OLI/TIRS cloud-free scenes between 2014 and 2023, covering major wet–dry cycles in Mauritania. With regard to this, the current paper contributed to mapping of the hydrologically vulnerable areas of Mauritanian coasts through detecting the fluctuating area of sebkhas and crust salt pan. Changes in the land cover types detected and visualised by GRASS GIS are generated from the Landsat 8-9 dataset and shown in the corresponding figures (Figures \ref{fig06}--\ref{fig11} according to the applied methods of unsupervised and supervised classification. RS data were used since they provide reliable source of information to clarify environmental dynamics of salt pans. In addition, they support further investigations on crucial limits on the sedimentological, mineralogical and hydrological history of the salt pans.

\newpar The land categories for years 2014–2013 and their statistics are presented in Table \ref{tab3}. Here, land cover classes are designated as the following categories: 1) Water bodies, 2) Shelf and coastal plains, 3) Sebkha, 4) Urban areas, 5) Sahelian grassland, 6) Salty sands, 7) Compact soil, 8) Stony desert and yellow dunes, 9) Sandy desert and white dunes and 10) Bare soil and rocks.

\begin{table}[H]
\caption{Estimated classes of land cover types in western Mauritania, Sebkha de Ndrhamcha, for April. Map units in measurements: 30 m resolution for each pixel on the multispectral scene of Landsat 8-9 OLI/TIRS}
  \centering
  \begin{tabularx}{\textwidth}{p{0.9cm}p{0.9cm}p{0.9cm}p{0.9cm}p{0.9cm}p{0.9cm}p{0.9cm}p{0.9cm}p{0.9cm}p{0.9cm}p{0.9cm}}
  \toprule
  \multirow{2}{*}{Year} & \multicolumn{10}{c}{Classes of land cover types in western Mauritania, Sebkha de Ndrhamcha} \\
     & 1 & 2 & 3 & 4 & 5 & 6 & 7 & 8 & 9 & 10 \\
    \midrule
    \textbf{2014} & 1410 & 4 & 30 & 79 & 202 & 708 & 1239 & 1581 & 1493 & 178 \\
    
    \textbf{2017} & 1386 & 43 & 70 & 141 & 215 & 668 & 1056 & 1742 & 1548 & 144 \\
    \textbf{2018} & 1384 & 21 & 78 & 110 & 182 & 598 & 1195 & 1975 & 1353 & 117 \\
    \textbf{2020} & 1395 & 2 & 28 & 145 & 228 & 582 & 881 & 2502 & 1231 & 20 \\
    \textbf{2022} & 1443 & 2 & 20 & 116 & 197 & 532 & 970 & 2567 & 1146 & 40 \\
    \textbf{2023} & 1468 & 31 & 53 & 136 & 253 & 538 & 1153 & 1622 & 1577 & 206 \\
\bottomrule
\end{tabularx}
  \label{tab3}
\end{table}

\newpar The results of the satellite image processing are presented in the maps shown in Figures 6 \ref{fig06}--\ref{fig11} in the following subsections. The results of the ML-based classification for Sebkha de Ndrhamcha region are shown in the corresponding figures for each method applied (Clustering, RF Classifier, DTC, GBC and SVM), and numerical computation of the area occupied by each land category is summarised in Table \ref{tab3}. The accuracy assessment of the classification model is evaluated using statistical algorithms for computation of the overall accuracy, User’s and Producer’s accuracy and Kappa coefficient. The results are reported in Table \ref{tab04}.

\newpar Here, the accuracy is evaluated as the total number of correctly classified pixels on the Landsat scenes divided by the total number of pixels on this image using User’s accuracy and Producer’s accuracy followed by Cohen’s Kappa computation. Hence, the User’s accuracy is defined as the proportion of pixels classified to a specific land cover type among all the evaluated pixels that really belong to that land cover type. The evaluation was performed against the ground truth data referred to the FAO classification and includes the computed false negatives. Table \ref{tab04} summarises and compares computed parameters of accuracy and Cohen’s Kappa coefficient.

\newpar The interpretation of the Cohen’s Kappa statistics is as follows. The highest possible value of 1.0 indicates the perfect precision and accurate classification of variables, while the lowest theoretically possible value of near zero indicates the lowest precision. For both User’s accuracy and Producer’s accuracy, higher values normally mean better results of the classification of the satellite images. Based on the computed values of Cohen’s Kappa coefficient and the accuracy method for the evaluation of classified data, the results suggest that the SVM and RF algorithms are rated as the most reliable approaches due to the high values.

\begin{longtable}{@{\extracolsep{\fill}}p{0.8cm}|p{0.8cm}p{0.8cm}p{0.8cm}p{0.8cm}|p{0.8cm}p{0.8cm}p{0.8cm}p{0.8cm}|p{0.8cm}p{0.8cm}p{0.8cm}p{0.8cm}@{}}
\caption{Accuracy assessment for ML models in GRASS GIS: 1) Random Forest (RF); 2) Support Vector Machine (SVM); 3) Decision Tree Classifier (DTC); 4) Gradient Boosting Classifier (GBC). Estimated classes of land cover types for 2014--2023 in West Mauritania.} \label{tab04}\\
\multirow{2}{*}{Year} & \multicolumn{4}{c}{Producer's accuracy, \%} & \multicolumn{4}{c}{User's accuracy, \%} & \multicolumn{4}{c}{Kappa statistics}\\
      & RF & SVM & DTC & GBC & RF & SVM & DTC & GBC & RF & SVM & DTC & GBC \\
 \endfirsthead
 
 \multicolumn{13}{c}%
{{\tablename\ \thetable{} -- continued from previous page}} \\
\multirow{2}{*}{Year} & \multicolumn{4}{c}{Producer's accuracy, \%} & \multicolumn{4}{c}{User's accuracy, \%} & \multicolumn{4}{c}{Kappa statistics}\\
      & RF & SVM & DTC & GBC & RF & SVM & DTC & GBC & RF & SVM & DTC & GBC \\ \hline
\endhead

\hline \multicolumn{13}{r}{{Continued on next page}} \\ \hline
\endfoot

\hline \hline
\endlastfoot

 \hline
    \multicolumn{13}{c}{Land Cover Class 1: Water bodies} \\
    \textbf{2014} & 77 & 83 & 74  & 71 & 78 & 88 & 62 & 62 & 0.80 & 0.74 & 0.65 & 0.65 \\
    \textbf{2017} & 78 & 84 & 72  & 70 & 77 & 85 & 65 & 67 & 0.81 & 0.93 & 0.53 & 0.59 \\
    \textbf{2018} & 76 & 79 & 75 & 69 & 71 & 79 & 66 & 68 & 1.00 & 0.90 & 0.69 & 0.61 \\
    \textbf{2020} & 81 & 81 & 73  & 72 & 89 & 89 & 68 & 72 & 0.93 & 0.95 & 0.58 & 0.61 \\
     \textbf{2022} & 82 & 80 & 65  & 73 & 82 & 73 & 75 & 74 & 0.95 & 0.70 & 0.84 & 1.00 \\
    \textbf{2023} & 79 & 82 & 66 & 68 & 77 & 74 & 74 & 70 & 0.74 & 0.79 & 0.83 & 0.67 \\
\hline
\multicolumn{13}{c}{Land Cover Class 2: Shelf and coastal plains} \\
    \textbf{2014} & 89 & 78 & 72  & 67 & 100 & 98 & 56 & 76 & 0.82 & 0.91 & 0.67 & 0.69 \\
    \textbf{2017} & 95 & 91 & 68  & 72 & 92 & 65 & 64 & 65 & 0.73 & 1.00 & 0.56 & 0.74 \\
    \textbf{2018} & 100 & 92 & 87  & 91 & 87 & 77 & 63 & 68 & 0.78 & 0.94 & 0.88 & 0.68 \\
    \textbf{2020} & 88 & 100 & 86  & 65 & 89 & 87 & 69 & 62 & 0.89 & 0.84 & 0.83 & 0.85 \\
     \textbf{2022} & 100 & 95 & 96  & 59 & 78 & 77 & 58 & 59 & 0.91 & 0.78 & 0.72 & 0.73 \\
    \textbf{2023} & 84 & 83 & 75  & 71 & 73 & 75 & 74 & 72 & 0.92 & 0.75 & 0.75 & 0.56 \\
\hline
\multicolumn{13}{c}{Land Cover Class 3: Sebkha} \\
    \textbf{2014} & 95 & 78 & 67  & 74 & 92 & 90 & 66 & 67 & 0.88 & 0.74 & 0.68 & 0.72 \\
    \textbf{2017} & 85 & 93 & 88  & 72 & 78 & 87 & 69 & 81 & 0.78 & 0.81 & 0.55 & 0.56 \\
    \textbf{2018} & 88 & 89 & 63  & 64 & 77 & 94 & 71 & 69 & 0.91 & 0.88 & 0.71 & 0.54 \\
    \textbf{2020} & 94 & 75 & 64  & 69 & 93 & 85 & 72 & 65 & 1.00 & 0.64 & 0.75 & 0.78 \\
     \textbf{2022} & 73 & 92 & 64  & 81 & 73 & 95 & 59 & 71 & 0.84 & 1.00 & 0.73 & 0.78 \\
    \textbf{2023} & 69 & 68 & 76  & 83 & 67 & 89 & 73 & 72 & 0.91 & 0.90 & 0.74 & 0.84 \\
\hline
\multicolumn{13}{c}{Land Cover Class 4: Urban areas} \\
    \textbf{2014} & 87 & 78 & 64  & 68 & 88 & 86 & 67 & 66 & 1.00 & 0.85 & 0.76 & 0.81 \\
    \textbf{2017} & 94 & 79 & 69  & 64 & 92 & 75 & 61 & 72 & 0.95 & 1.00 & 0.82 & 0.73 \\
    \textbf{2018} & 74 & 91 & 74  & 72 & 91 & 79 & 72 & 74 & 0.85 & 0.87 & 0.75 & 0.80 \\
    \textbf{2020} & 69 & 80 & 65  & 76 & 84 & 64 & 68 & 68 & 1.00 & 0.79 & 0.81 & 0.69 \\
     \textbf{2022} & 71 & 66 & 83  & 81 & 74 & 62 & 54 & 65 & 0.92 & 1.00 & 0.65 & 0.75 \\
    \textbf{2023} & 83 & 65 & 72  & 83 & 79 & 75 & 59 & 71 & 0.87 & 0.72 & 0.69 & 0.73 \\
\hline
\multicolumn{13}{c}{Land Cover Class 5: Sahelian grassland} \\
    \textbf{2014} & 87 & 81 & 65  & 67 & 82 & 90 & 83 & 81 & 0.93 & 0.74 & 0.95 & 0.91 \\
    \textbf{2017} & 64 & 65 & 71  & 68 & 76 & 79 & 74 & 78 & 0.78 & 0.68 & 0.76 & 0.58 \\
    \textbf{2018} & 78 & 72 & 73  & 57 & 79 & 91 & 68 & 63 & 1.00 & 0.83 & 0.73 & 0.79 \\
    \textbf{2020} & 74 & 79 & 72  & 71 & 67 & 77 & 64 & 68 & 0.91 & 0.65 & 0.88 & 0.63 \\
     \textbf{2022} & 88 & 82 & 71  & 82 & 81 & 80 & 72 & 69 & 0.53 & 1.00 & 0.54 & 0.77 \\
    \textbf{2023} & 91 & 75 & 65  & 64 & 91 & 82 & 71 & 90 & 1.00 & 0.91& 0.67 & 0.82 \\
\hline
\multicolumn{13}{c}{Land Cover Class 6: Salty sands} \\
    \textbf{2014} & 88 & 72 & 67  & 63 & 91 & 79 & 68 & 67 & 0.98 & 0.73 & 0.73 & 0.78 \\
    \textbf{2017} & 72 & 79 & 72  & 76 & 82 & 74 & 65 & 78 & 1.00 & 0.75 & 0.81 & 0.75 \\
    \textbf{2018} & 91 & 90 & 64  & 78 & 90 & 78 & 71 & 81 & 1.00 & 0.81 & 0.76 & 0.69\\
    \textbf{2020} & 89 & 85 & 69  & 74 & 85 & 81 & 78 & 82 & 0.81 & 0.86 & 0.88 & 0.65 \\
     \textbf{2022} & 90 & 84 & 56  & 71 & 78 & 90 & 64 & 89 & 0.84 & 0.91 & 0.72 & 0.67 \\
    \textbf{2023} & 77 & 71 & 68  & 80 & 74 & 72 & 73 & 75 & 0.78 & 1.00 & 0.63 & 0.68 \\
\hline
\multicolumn{13}{c}{Land Cover Class 7: Compact soil} \\
    \textbf{2014} & 89 & 78 & 67 & 63 & 89 & 81 & 69 & 63 & 0.88 & 0.73 & 0.65 & 0.61 \\
    \textbf{2017} & 73 & 77 & 63 & 69 & 73 & 72 & 73 & 71 & 1.00 & 0.61 & 0.69 & 0.83 \\
    \textbf{2018} & 78 & 69 & 72 & 65 & 85 & 85 & 74 & 70 & 0.85 & 0.68 & 0.81 & 0.95 \\
    \textbf{2020} & 88 & 81 & 79 & 71 & 69 & 74 & 65 & 65 & 0.97 & 0.74 & 0.85 & 0.74 \\
     \textbf{2022} & 91 & 90 & 80 & 75 & 70 & 78 & 69 & 69 & 0.83 & 1.00 & 0.72 & 0.61 \\
    \textbf{2023} & 74 & 75 & 71 & 61 & 66 & 64 & 81 & 74 & 0.74 & 0.92 & 0.78 & 0.89 \\
\hline
\multicolumn{13}{c}{Land Cover Class 8: Stony desert and yellow dunes} \\
    \textbf{2014} & 85 & 83 & 75  & 67 & 89 & 73 & 68 & 74 & 0.84 & 0.81 & 0.71 & 0.86 \\
    \textbf{2017} & 78 & 88 & 78  & 69 & 93 & 84 & 69 & 82 & 0.75 & 1.00 & 0.82 & 0.92 \\
    \textbf{2018} & 89 & 81 & 74  & 78 & 78 & 86 & 89 & 76 & 1.00 & 0.63 & 0.70 & 0.84 \\
    \textbf{2020} & 92 & 79 & 91  & 71 & 95 & 72 & 81 & 81 & 0.72 & 0.79 & 0.62 & 0.70 \\
     \textbf{2022} & 91 & 84 & 73  & 74 & 81 & 89 & 90 & 89 & 0.68 & 0.71 & 0.59 & 0.69 \\
    \textbf{2023} & 84 & 75 & 75  & 73 & 73 & 90 & 83 & 73 & 0.66 & 0.65 & 0.60 & 0.76 \\
\hline
\multicolumn{13}{c}{Land Cover Class 9: Sandy desert and white dunes} \\
    \textbf{2014} & 92 & 81 & 67  & 65 & 81 & 82 & 65 & 71 & 0.77 & 0.82 & 0.77 & 0.71 \\
    \textbf{2017} & 83 & 89 & 72  & 81 & 84 & 78 & 78 & 78 & 0.81 & 0.73 & 0.81 & 0.55 \\
    \textbf{2018} & 93 & 70 & 68  & 72 & 95 & 73 & 71 & 73 & 0.64 & 0.67 & 0.80 & 1.00 \\
    \textbf{2020} & 78 & 93 & 80  & 68 & 88 & 89 & 83 & 67 & 1.00 & 0.54 & 0.68 & 0.75 \\
     \textbf{2022} & 77 & 94 & 71  & 73 & 89 & 71 & 69 & 81 & 0.98 & 0.91 & 0.72 & 0.67 \\
    \textbf{2023} & 74 & 85 & 81  & 79 & 73 & 69 & 62 & 66 & 0.90 & 1.00 & 0.74 & 0.79 \\
\hline
\multicolumn{13}{c}{Land Cover Class 10: Bare soil and rocks} \\
    \textbf{2014} & 91 & 83 & 68  & 78 & 84 & 83 & 78 & 74 & 0.87 & 1.00 & 0.73 & 0.69 \\
    \textbf{2017} & 89 & 82 & 63  & 72 & 83 & 78 & 73 & 75 & 0.81 & 0.74 & 0.77 & 0.57 \\
    \textbf{2018} & 88 & 75 & 65  & 77 & 79 & 69 & 69 & 81 & 1.00 & 0.73 & 0.81 & 0.61 \\
    \textbf{2020} & 79 & 69 & 73  & 64 & 74 & 71 & 74 & 67 & 0.74 & 0.69 & 0.80 & 0.74 \\
     \textbf{2022} & 82 & 91 & 79  & 69 & 75 & 73 & 81 & 83 & 0.63 & 0.81 & 1.00 & 0.83 \\
    \textbf{2023} & 77 & 80 & 75  & 71 & 82 & 95 & 83 & 74 & 0.84 & 0.68 & 0.65 & 0.80 \\
    \hline
\end{longtable}

\newpar The Producer's accuracy indicates the proportion of the pixels that are correctly classified and belong to a target land cover type. The Cohen's Kappa coefficient measures the difference between the actual reference between the data based on the Overall accuracy and the actual results of the land cover classification. In other words, it shows the actual changes between the referenced data and a randomly selected pixel classified using given method. As other correlation statistics, the Kappa values can range from --1 to +1. The interpretation of values is as follows. Values of the levels above 60\% to 80\% have a moderate accuracy of pixel's classification, while those above 80\% are acceptable and generally demonstrate accurate identification of target class. 

\newpar Consequently, the values approaching above 90\% and nearly 100\% indicate almost perfect classification results. Thus, the SVM and RF techniques can be evaluated as a robust approaches of GRASS GIS for RS data classification based on accuracy, followed by the DTC and GBC methods. The overall classification accuracy is above 80\% with a corresponding overall Kappa coefficient above 0.85 (Table \ref{tab04}). Such results demonstrate robust conceptual agreement of the processed data with ground truth. This indicates that GRASS GIS has a reliable algorithms of RS data processing that can be used in similar works.

\subsection*{Clustering}

\begin{figure}[H]
	\begin{subfigure}[b]{.5\textwidth}
		\centering
			\includegraphics[width=0.80\textwidth]{Fig_06a}
		\caption{2014}
	\end{subfigure}%
	\begin{subfigure}[b]{.5\textwidth}
		\centering
		\includegraphics[width=0.81\textwidth]{Fig_06b}
		\caption{2017}
	\end{subfigure}%
\\
\vfill \vspace{1mm}
	\begin{subfigure}[b]{.5\textwidth}
		\centering
			\includegraphics[width=0.80\textwidth]{Fig_06c}
		\caption{2018}
	\end{subfigure}%
	\begin{subfigure}[b]{.5\textwidth}
		\centering
		\includegraphics[width=0.81\textwidth]{Fig_06d}
		\caption{2020}
	\end{subfigure}%
\\
\vfill \vspace{1mm}
	\begin{subfigure}[b]{.5\textwidth}
		\centering
			\includegraphics[width=0.80\textwidth]{Fig_06e}
		\caption{2022}
	\end{subfigure}%
	\begin{subfigure}[b]{.5\textwidth}
		\centering
		\includegraphics[width=0.81\textwidth]{Fig_06f}
		\caption{2023}
	\end{subfigure}%
\vspace*{20pt}\caption{Classified Landsat 8-9 OLI/TIRS images using clustering: (\textbf{a}) 2014, (\textbf{b}) 2017, (\textbf{c}) 2018, (\textbf{d}) 2020, (\textbf{e}) 2022, (\textbf{f}) 2023. Software: GRASS GIS. Source of maps: author.}\label{fig06}
\end{figure}

\newpar The accuracy assessment performed using the estimated rejection probability is shown in Figure \ref{fig07}. It shows the correctness of the pixels associated with each land cover class according to the clustering method of k-means, used in GRASS GIS. The map of land cover types for 2014 had an overall kappa statistic of 0.97 and an overall accuracy of 98.34\% which indicated the correctly classified pixels with stable points, the classification for 2017 had overall accuracy of 98.08\%, for year 2018 -- 98.29\%, for 2020 -- 98.36\%, for 2022 -- 98.64\% and for year 2023 -- 98.56\% of computed overall accuracy.  

\begin{figure}[H]
	\begin{subfigure}[b]{.5\textwidth}
		\centering
			\includegraphics[width=0.80\textwidth]{Fig_07a}
		\caption{2014}
	\end{subfigure}%
	\begin{subfigure}[b]{.5\textwidth}
		\centering
		\includegraphics[width=0.80\textwidth]{Fig_07b}
		\caption{2017}
	\end{subfigure}%
\\
\vfill \vspace{1mm}
	\begin{subfigure}[b]{.5\textwidth}
		\centering
			\includegraphics[width=0.80\textwidth]{Fig_07c}
		\caption{2018}
	\end{subfigure}%
	\begin{subfigure}[b]{.5\textwidth}
		\centering
		\includegraphics[width=0.80\textwidth]{Fig_07d}
		\caption{2020}
	\end{subfigure}%
\\
\vfill \vspace{1mm}
	\begin{subfigure}[b]{.5\textwidth}
		\centering
			\includegraphics[width=0.80\textwidth]{Fig_07e}
		\caption{2022}
	\end{subfigure}%
	\begin{subfigure}[b]{.5\textwidth}
		\centering
		\includegraphics[width=0.80\textwidth]{Fig_07f}
		\caption{2023}
	\end{subfigure}%
\vspace*{20pt}\caption{Maps of rejection threshold probability for accuracy analysis of image classification by chi-squared test: (\textbf{a}) 2014, (\textbf{b}) 2017, (\textbf{c}) 2018, (\textbf{d}) 2020, (\textbf{e}) 2022, (\textbf{f}) 2023. Software: GRASS GIS. Source of maps: author.}\label{fig07}
\end{figure}

%---------------------ML -------------->

\subsection*{Random Forest Classifier}

The results of the Random Forest classification applied for mapping sebkhas are shown in Figure \ref{fig08}. The following 10 land cover classes were identified and visualised, Figure \ref{fig06}: 1) Water bodies; 2) Shelf and coastal plains; 3) Sebkha; 4) Urban areas; 5) Sahelian grassland; 6) Salty sands; 7) Compact soil; 8) Stony desert and yellow dunes; 9) Sandy desert and white dunes; 10) Bare soil and rocks. 

\begin{figure}[H]
	\begin{subfigure}[b]{.5\textwidth}
		\centering
			\includegraphics[width=0.80\textwidth]{Fig_08a}
		\caption{2014}
	\end{subfigure}%
	\begin{subfigure}[b]{.5\textwidth}
		\centering
		\includegraphics[width=0.80\textwidth]{Fig_08b}
		\caption{2017}
	\end{subfigure}%
\\
\vfill \vspace{1mm}
	\begin{subfigure}[b]{.5\textwidth}
		\centering
			\includegraphics[width=0.80\textwidth]{Fig_08c}
		\caption{2018}
	\end{subfigure}%
	\begin{subfigure}[b]{.5\textwidth}
		\centering
		\includegraphics[width=0.80\textwidth]{Fig_08d}
		\caption{2020}
	\end{subfigure}%
\\
\vfill \vspace{1mm}
	\begin{subfigure}[b]{.5\textwidth}
		\centering
			\includegraphics[width=0.80\textwidth]{Fig_08e}
		\caption{2022}
	\end{subfigure}%
	\begin{subfigure}[b]{.5\textwidth}
		\centering
		\includegraphics[width=0.81\textwidth]{Fig_08f}
		\caption{2023}
	\end{subfigure}%
\vspace*{20pt}\caption{Random Forest Classifier of Machine Learning (ML) classification methods applied for processing of satellite images: (\textbf{a}) 2014, (\textbf{b}) 2017, (\textbf{c}) 2018, (\textbf{d}) 2020, (\textbf{e}) 2022, (\textbf{f}) 2023. Software: GRASS GIS. Source of maps: author.}\label{fig08}
\end{figure}

\subsection*{Decision Tree Classifier}

Based on the presented results, the largest category was barren land subdivided into different classes (Salty sands, Compact soil, Stony desert and yellow dunes, Sandy desert and white dunes), followed by vegetation class (Sahelian grassland) and salt pan (Class 2 of sebkha). The other land use categories included water surfaces, built-up urban areas and agricultural lands. Using the maps demonstrated on the series of images processed by Decision Tree classification, the distinct land cover classes are visualised for the study area, including the location of sebkhas and its changing extent over the years (Figure \ref{fig09}).

\begin{figure}[H]
	\begin{subfigure}[b]{.5\textwidth}
		\centering
			\includegraphics[width=0.80\textwidth]{Fig_09a}
		\caption{2014}
	\end{subfigure}%
	\begin{subfigure}[b]{.5\textwidth}
		\centering
		\includegraphics[width=0.80\textwidth]{Fig_09b}
		\caption{2017}
	\end{subfigure}%
\\
\vfill \vspace{1mm}
	\begin{subfigure}[b]{.5\textwidth}
		\centering
			\includegraphics[width=0.80\textwidth]{Fig_09c}
		\caption{2018}
	\end{subfigure}%
	\begin{subfigure}[b]{.5\textwidth}
		\centering
		\includegraphics[width=0.80\textwidth]{Fig_09d}
		\caption{2020}
	\end{subfigure}%
\\
\vfill \vspace{1mm}
	\begin{subfigure}[b]{.5\textwidth}
		\centering
			\includegraphics[width=0.80\textwidth]{Fig_09e}
		\caption{2022}
	\end{subfigure}%
	\begin{subfigure}[b]{.5\textwidth}
		\centering
		\includegraphics[width=0.80\textwidth]{Fig_09f}
		\caption{2023}
	\end{subfigure}%
\vspace*{20pt}\caption{Decision Tree Classifier of Machine Learning (ML) classification methods applied for processing of satellite images: (\textbf{a}) 2014, (\textbf{b}) 2017, (\textbf{c}) 2018, (\textbf{d}) 2020, (\textbf{e}) 2022, (\textbf{f}) 2023. Software: GRASS GIS. Source of maps: author.}\label{fig09}
\end{figure}

\subsection*{Gradient Boosting Classifier}

The areas of diverse land cover types detected using Gradient Boosting classification method and their changes over the evaluated periods are presented in Figure \ref{fig10}. According to the analysis on satellite images processed by SVMs, the largest extent of wetlands was reached in 2017 and it reduced in consecutive years (2018–2022), after which the extent recovered in 2023 (Figure \ref{fig10}).

\begin{figure}[H]
	\begin{subfigure}[b]{.5\textwidth}
		\centering
			\includegraphics[width=0.80\textwidth]{Fig_10a}
		\caption{2014}
	\end{subfigure}%
	\begin{subfigure}[b]{.5\textwidth}
		\centering
		\includegraphics[width=0.81\textwidth]{Fig_10b}
		\caption{2017}
	\end{subfigure}%
\\
\vfill \vspace{1mm}
	\begin{subfigure}[b]{.5\textwidth}
		\centering
			\includegraphics[width=0.80\textwidth]{Fig_10c}
		\caption{2018}
	\end{subfigure}%
	\begin{subfigure}[b]{.5\textwidth}
		\centering
		\includegraphics[width=0.81\textwidth]{Fig_10d}
		\caption{2020}
	\end{subfigure}%
\\
\vfill \vspace{1mm}
	\begin{subfigure}[b]{.5\textwidth}
		\centering
			\includegraphics[width=0.80\textwidth]{Fig_10e}
		\caption{2022}
	\end{subfigure}%
	\begin{subfigure}[b]{.5\textwidth}
		\centering
		\includegraphics[width=0.81\textwidth]{Fig_10f}
		\caption{2023}
	\end{subfigure}%
\vspace*{20pt}\caption{Gradient Boosting Classifier of Machine Learning (ML) classification methods applied for processing of satellite images: (\textbf{a}) 2014, (\textbf{b}) 2017, (\textbf{c}) 2018, (\textbf{d}) 2020, (\textbf{e}) 2022, (\textbf{f}) 2023. Software: GRASS GIS. Source of maps: author.}\label{fig10}
\end{figure}

\subsection*{SVM Classifier}

The SVM method of ML approach of GRASS GIS enabled to reveal hidden information on land cover types and changes in landscape patterns in West Mauritania through processed satellite Landsat images (Figure \ref{fig11}). 

\begin{figure}[H]
	\begin{subfigure}[b]{.5\textwidth}
		\centering
			\includegraphics[width=0.80\textwidth]{Fig_11a}
		\caption{2014}
	\end{subfigure}%
	\begin{subfigure}[b]{.5\textwidth}
		\centering
		\includegraphics[width=0.80\textwidth]{Fig_11b}
		\caption{2017}
	\end{subfigure}%
\\
\vfill \vspace{1mm}
	\begin{subfigure}[b]{.5\textwidth}
		\centering
			\includegraphics[width=0.80\textwidth]{Fig_11c}
		\caption{2018}
	\end{subfigure}%
	\begin{subfigure}[b]{.5\textwidth}
		\centering
		\includegraphics[width=0.81\textwidth]{Fig_11d}
		\caption{2020}
	\end{subfigure}%
\\
\vfill \vspace{1mm}
	\begin{subfigure}[b]{.5\textwidth}
		\centering
			\includegraphics[width=0.80\textwidth]{Fig_11e}
		\caption{2022}
	\end{subfigure}%
	\begin{subfigure}[b]{.5\textwidth}
		\centering
		\includegraphics[width=0.80\textwidth]{Fig_11f}
		\caption{2023}
	\end{subfigure}%
\vspace*{20pt}\caption{Support Vector Machine (SVM) Classifier of Machine Learning (ML) classification methods applied for processing of satellite images: (\textbf{a}) 2014, (\textbf{b}) 2017, (\textbf{c}) 2018, (\textbf{d}) 2020, (\textbf{e}) 2022, (\textbf{f}) 2023. Software: GRASS GIS. Source of maps: author.}\label{fig11}
\end{figure}

\newpar The results revealed that the largest land cover class was Class 8 (Stony desert and yellow dunes) according to the extent of the satellite images, followed by Class 7 (Compact soil and barren land) in the desert areas, which represent the dunes of the Sahara. It was then followed by water surfaces (Class 1), which occupy shelf areas, Class 7 (Compact soil of the pre-Sahara lands) and Class 6 (Salty sands located in the saline soils). The distribution of other land cover types is summarised in Table 3. The smallest class is occupied by Class 2, which is represented by the narrow shelf strip of the coastal plains with specific aquatic vegetation between the water area and coastal regions.

\newpar It is worth comparing the classifier algorithms used in this study—RF Classifier, DTC, GBC and SVM Classifier—with CNNs, which provides the following remarks. In contrast to the ML methods, Deep Learning (DL) approaches have also increasingly been applied to the tasks of spatial data processing, including image classification or segmentation. Among them, the CNN method has distinct methodological framework. Thus, in contrast to RF or SVM, CNN consists of a sequence of three layers: convolution, non-linearity and pooling.

\newpar Repeated iteratively several times as gradient descent, this operation is followed by a step to determine the land cover class with the highest probability using the final fully connected layer. In this way, CNN discovers the most discriminative features in a satellite image and evaluates the context of the surrounding landscapes using weights of the Neural Network (NN) learnt through previous iterative stochastic algorithm. Hence, in comparison to SVM, CNN is distinct of being non-parametric and able to handle large volumes of data with low risk of overfitting. Such properties of CNN are useful for bid data processing, for example, large series of images. Besides, CNN is generally quick and easy to configure using programming tools.

\newpar The ML used for image classification in this study demonstrated robust results for diverse tasks of image processing. Some of them are land cover classification, recognition of buildings and geometric objects on the images and classification of diverse land surface types to evaluate environmental dynamics. For instance, the thematic maps can be used for visualising the landscapes affected by tsunami or floods, which can be detected using times series of images before and after the catastrophic events. Besides, the ML approaches work well in object detection tasks. As demonstrated in this study, the accuracy of the land cover classifications using ML methods demonstrated the highest results for RF and SVMs, which reflects the ability of these methods to exploit, through embedded algorithms, spatial and textural features on the processed images. Moreover, these methods enable extracting spectral information from the spatial patterns on the satellite imagery. The ability of ML methods to process this information suggests their application in similar studies that use moderate to resolution imagery for environmental monitoring and mapping.
 
\newsection {4. CONCLUSIONS}

\newpar This paper showed the application of GRASS GIS modules of ML for satellite image processing. Practical use of the presented maps consists in the support to make use of geoinformation, adjust it for land monitoring and discover the climate factors which cause landscape dynamics in the coastal regions of Africa. Using the algorithms of image processing from statistical and programming libraries, the changed landscape patterns were detected and the cartographic visualisation was adjusted from a computer vision perspective using several ML-based algorithms, which are RF Classifier, DTC, GBC and SVM Classifier. The traditional unsupervised methods of classification were compared with the ML-based methods. While clustering demonstrated difficulties in correct detection of land parcels using spectral reflectance, the ML-based supervised techniques of trained learning, as a complement, detected, identified and recorded such information using computer vision algorithms in an automated manner. The outcomes can also be employed to enhance or control the quality of current land cover maps in Mauritania.

\newpar To cope with the problem of soil salinisation in the drought-prone region of arid Africa, such as coastal Mauritania, EO data are needed as effective data source for comparison and visualisation of changes and dynamics of salt crust. However, the use of the traditional methods of RS data processing is time-consuming and might be prone to misclassification. In contrast, ML methods applied to EO data processing ensure automated and effective way to visualise, classify and numerically estimate changes in land cover types using algorithms of statistical data processing and computer vision. With regard to this, this manuscript  contributes to the development of such methods in West Africa.

\newpar This study presents the ML-based cartographic analysis of changes in land cover changes in the coastal region of southwestern Mauritania, including the salt pan of Sebkha de Ndrhamcha, to visualise dynamics in salt crust and surrounding landscapes. The data and methods were selected to achieve this goal according to the cartographic techniques of satellite image processing. Hence, the data used for evaluation included satellite Landsat images obtained from USGS as the RS dataset from 2014 to 2023. The methods used for evaluation and cartographic work included unsupervised and supervised classification used for comparison between traditional and ML GIS techniques of GRASS GIS. The results indicated that there was a change in land cover types. The changes between the class occupied by compact soil and bare soil and rocks correlate well with those in sandy desert and dunes. This is especially notable for the years 2020 and 2022, which shows regional climate impacts with the effects of winds and aeolian processes that increase the areas of dunes and sandy areas and decrease bare soils, accordingly. Likewise, the decrease of shelf and coastal plains is visible for the same period, which might be related to the active upwelling processes in this time.

\newpar In general, the coastal region of southwestern Mauritania, along Sebkha de Ndrhamcha, has experienced considerable changes in land cover classes over the past decade (2014–2023), throughout the area of categories, including their fragmentation and extent. However, changes have been significant in the central part of the salt pan, which also supports the main fluctuations of saline crust and salinisation of soils in the surroundings. Thus, the results of this study revealed that over the past 10-year period, salinisation of the sandy areas around salt pan increased, whereas the barren lands decreased in the study area and were replaced by the salty sands and dunes of desert. The details of land cover changes are shown in corresponding figures made with various techniques of ML algorithms using GRASS GIS. The overall results of the study highlighted that GRASS GIS techniques of ML are highly effective for satellite image processing as advanced and straightforward methods of supervised image classification.

\newpar The satellite-based evaluation of landscape in western Mauritania is focused on the Sahara–Sahel transition zone of arid and semi-arid climate. Scarce water resources and fluctuation of salt lakes and sebhka motivate the need for mapping of water availability and distribution of salt pans. Changes in land cover classes including fluctuations in salinisation were revealed using ML methods of satellite image processing using GRASS GIS. Practical application of this work consists in land management since visualisation of the extents of wetlands and water bodies enables to plan potential drainage networks based on the hydrological connections, as well as evaluate spatial location of aquatic habitats \cite{CAMPOS2012438}. Hence, new created maps demonstrated changes in distribution of sebkhas and can assist land planners and environmentalists specialised on Mauritania.

\newpar The limited scope of this study was focused on the preselected tools with the available libraries of GRASS GIS, which present an interface to the scikit-learn library of Python. Therefore, future similar research can be proposed with the focus on further development of spatiotemporal analysis using GRASS GIS. For example, the use of GRASS GIS with the available ML tools may be extended towards its more sophisticated tools for spatiotemporal analysis (t.* commands). The quality control and enhancement of current land cover products in Mauritania could be achieved through land cover maps produced using such ML techniques. Furthermore, we also recommend the use of other satellite products such as Sentinel-2. The application of RS data with high resolution would be useful for broad-scale land cover mapping in West Africa over selected time periods. The Sentinel-2 data can be achieved based on Copernicus open-access repository. Finally, future work is suggested to produce and evaluate the effectiveness of a satellite time series of broad scale covering Mauritania for annual land cover mapping. The comparison of such data with existing reference land cover products would be useful for operative monitoring of arid areas of West Africa.

\newpar To conclude, the demonstrated case of environmental mapping by satellite data processing by using ML algorithms supports the analysis of salinisation in salty sand areas and sebkha of West Mauritania related to climate change. Using GRASS GIS as a powerful cartographic tool promotes the use of scripts as advanced techniques for processing geoinformation. Such technologies ensure future research efforts, which can continue further the presented study through focus on identifying salinisation of soil for other time periods. GRASS GIS-based methods of processing EO data and formulating optimal ML approaches among the existing algorithms enable to develop effective methods for satellite image processing. In turn, efficient cartographic instruments supported by programming methods for image processing provide technical support for the environmental climate visualisation and analysis of rare areas such as Sebkha de Ndrhamcha, West Mauritania.

\vspace{12pt}

\newsection {ABBREVIATIONS}
\vspace{12pt}

\begin{tabular}{@{}ll}
AI & Artificial Intelligence \\
CNN & Convolutional Neural Network \\
DCW & Digital Chart of the World \\
DL & Deep Learning \\
EO & Earth Observation\\
GEBCO & General Bathymetric Chart of the Oceans \\
GIS & Geographic Information System \\
GMT & Generic Mapping Tools \\
GRASS & Geographic Resources Analysis Support System \\
Landsat OLI/TIRS & Landsat Operational Land Imager and the Thermal Infrared Sensor \\
NASA & National Aeronautics and Space Administration \\
NN & Neural Network \\
SRTM & Shuttle Radar Topography Mission \\
USGS & United States Geological Survey
\end{tabular}

\vspace{12pt}

\newsection {APPENDIX}\label{App}
\vspace{12pt}

\subsection*{GRASS GIS scripts for RS data processing}

\begin{lstlisting}[language=bash, caption=GRASS GIS script for data import and preprocessing,basicstyle=\ttfamily\scriptsize,keywordstyle=\color{deepmagenta},emphstyle=\color{deepred},
stringstyle=\color{deepblue},breaklines=true,numbers=left,label={lst01}]
r.import input=/Users/polinalemenkova/grassdata/Mauritania/\
	LC09_L2SP_205047_20230428_20230430_02_T1_SR_B1.TIF 
	output=L_2023_01 extent=region resolution=region
r.import input=/Users/polinalemenkova/grassdata/Mauritania/
	LC09_L2SP_205047_20230428_20230430_02_T1_SR_B2.TIF 
	output=L_2023_02 extent=region resolution=region
r.import input=/Users/polinalemenkova/grassdata/Mauritania/
	LC09_L2SP_205047_20230428_20230430_02_T1_SR_B3.TIF 
	output=L_2023_03 extent=region resolution=region
r.import input=/Users/polinalemenkova/grassdata/Mauritania/
	LC09_L2SP_205047_20230428_20230430_02_T1_SR_B4.TIF 
	output=L_2023_04 extent=region resolution=region
r.import input=/Users/polinalemenkova/grassdata/Mauritania/
	LC09_L2SP_205047_20230428_20230430_02_T1_SR_B5.TIF 
	output=L_2023_05 extent=region resolution=region
r.import input=/Users/polinalemenkova/grassdata/Mauritania/
	LC09_L2SP_205047_20230428_20230430_02_T1_SR_B6.TIF 
	output=L_2023_06 extent=region resolution=region
r.import input=/Users/polinalemenkova/grassdata/Mauritania/
	LC09_L2SP_205047_20230428_20230430_02_T1_SR_B7.TIF 
	output=L_2023_07 extent=region resolution=region
# check up the raster files
g.list rast
# adding shaded relief
r.import input=/Users/polinalemenkova/grassdata/
	Mauritania/gebco_2023.tif output=shaded_relief extent=region --overwrite
r.contour shaded_relief out=isolines step=200 --overwrite
\end{lstlisting}

\begin{lstlisting}[language=bash, caption=GRASS GIS script for image clustering and classification,basicstyle=\ttfamily\scriptsize,keywordstyle=\color{deepmagenta},emphstyle=\color{deepred},
stringstyle=\color{deepblue},breaklines=true,numbers=left,label={lst02}]
# Setting computational region to match the scene
g.region raster=L_2023_01 -p
# grouping data by i.group
i.group group=L_2023 subgroup=res_30m \
  input=L_2023_01,L_2023_02,L_2023_03,L_2023_04,\
  	L_2023_05,L_2023_06,L_2023_07 --overwrite
# Generating signature file and report using k-means clustering algorithm
i.cluster group=L_2023 subgroup=res_30m \
  signaturefile=cluster_L_2023 \
  classes=10 reportfile=rep_clust_L_2023.txt --overwrite
# Classification by i.maxlik module
i.maxlik group=L_2023 subgroup=res_30m \
  signaturefile=cluster_L_2023 \
  output=L_2023_clusters reject=L_2023_cluster_reject --overwrite
r.colors L_2023_clusters color=viridis
r.colors shaded_relief color=grey -n
# Mapping
g.region raster=L_2023_01 -p
d.mon wx0
d.rast shaded_relief
d.rast L_2023_clusters
d.vect isolines color='100:93:134' width=0
d.grid -g size=00:30:00 color=white width=0.1 fontsize=16 text_color=white
d.legend raster=L_2023_clusters title="Clusters 2023" title_fontsize=19 \
	font="Helvetica" fontsize=17 bgcolor=white border_color=white
d.legend raster=shaded_relief title="Relief, m" title_fontsize=19 \
	font="Helvetica" fontsize=17 bgcolor=white border_color=white -f
d.out.file output=Mauritania_2023 format=jpg --overwrite
# Mapping rejection probability
d.mon wx2
g.region raster=L_2023_clusters -p
r.colors L_2023_cluster_reject color=rainbow -e
d.rast shaded_relief
d.rast L_2023_cluster_reject
d.vect isolines color='100:93:134' width=0
d.grid -g size=00:30:00 color=white width=0.1 fontsize=16 text_color=white
d.legend raster=L_2023_cluster_reject title="2023" title_fontsize=19 \
	font="Helvetica" fontsize=17 bgcolor=white border_color=white
d.legend raster=shaded_relief title="Relief, m" title_fontsize=19 \
	font="Helvetica" fontsize=17 bgcolor=white border_color=white -f
d.out.file output=Mauritania_2023_reject format=jpg --overwrite
\end{lstlisting}

\subsection*{Machine Learning scripts of GRASS GIS for image processing}

\begin{lstlisting}[language=bash, caption=Random Forest Classifier script for image classification,basicstyle=\ttfamily\scriptsize,keywordstyle=\color{deepmagenta},emphstyle=\color{deepred},
stringstyle=\color{deepblue},breaklines=true,numbers=left,label={lst03}]
# Generating training pixels from an older land cover classification:
r.random input=L_2014_clusters seed=100 npoints=1000 raster=training_pixels
# training a model using r.learn.train
r.learn.train group=L_2023 training_map=training_pixels \
    model_name=RandomForestClassifier n_estimators=500 
    save_model=rfc_model.gz --overwrite
# perform prediction using r.learn.predict
r.learn.predict group=L_2023 load_model=rfc_model.gz 
output=rfc_classification_2023 --overwrite
# check raster categories which are applied to the classification output
r.category rfc_classification_2023
# display the results
r.colors rfc_classification_2023 color=roygbiv -e
d.mon wx0
d.rast shaded_relief
d.vect isolines color='100:93:134' width=0
d.rast rfc_classification_2023
d.grid -g size=00:30:00 color=white width=0.1 fontsize=16 text_color=white
d.legend raster=rfc_classification_2023 title="RFC 2023" title_fontsize=19 
	font="Helvetica" fontsize=17 bgcolor=white border_color=white
d.legend raster=shaded_relief title="Relief, m" title_fontsize=19 
	font="Helvetica" fontsize=17 bgcolor=white border_color=white -f
d.out.file output=RFC_2023 format=jpg --overwrite
\end{lstlisting}

\begin{lstlisting}[language=bash, caption=Decision Tree Classifier script for image classification,basicstyle=\ttfamily\scriptsize,keywordstyle=\color{deepmagenta},emphstyle=\color{deepred},
stringstyle=\color{deepblue},breaklines=true,numbers=left,label={lst04}]
# training a model using r.learn.train
r.learn.train group=L_2023 training_map=training_pixels \
    model_name=DecisionTreeClassifier n_estimators=500 
    	save_model=dtc_model.gz --overwrite
# perform prediction using r.learn.predict
r.learn.predict group=L_2023 load_model=dtc_model.gz 
	output=dtc_classification_2023 --overwrite
# check raster categories automatically applied to the classification output
r.category dtc_classification_2023
# display
r.colors dtc_classification_2023 color=bcyr
d.mon wx0
d.rast shaded_relief
d.vect isolines color='100:93:134' width=0
d.rast dtc_classification_2023
d.grid -g size=00:30:00 color=white width=0.1 fontsize=16 text_color=white
d.legend raster=dtc_classification_2023 title="DTC 2023" title_fontsize=19 
	font="Helvetica" fontsize=17 bgcolor=white border_color=white
d.legend raster=shaded_relief title="Relief, m" title_fontsize=19 
	font="Helvetica" fontsize=17 bgcolor=white border_color=white -f
d.out.file output=DTC_2023 format=jpg --overwrite
\end{lstlisting}

\begin{lstlisting}[language=bash, caption=Gradient Boosting Classifier script for image classification,basicstyle=\ttfamily\scriptsize,keywordstyle=\color{deepmagenta},emphstyle=\color{deepred},
stringstyle=\color{deepblue},breaklines=true,numbers=left,label={lst05}]
# training a model using r.learn.train
r.learn.train group=L_2023 training_map=training_pixels \
    model_name=GradientBoostingClassifier n_estimators=500 
    	save_model=gbc_model.gz --overwrite
# perform prediction using r.learn.predict
r.learn.predict group=L_2023 load_model=gbc_model.gz 
	output=gbc_classification_2023 --overwrite
# check raster categories automatically applied to the classification output
r.category gbc_classification_2023
# display
r.colors gbc_classification_2023 color=bgyr
d.mon wx0
d.rast shaded_relief
d.vect isolines color='100:93:134' width=0
d.rast gbc_classification_2023
d.grid -g size=00:30:00 color=white width=0.1 fontsize=16 text_color=white
d.legend raster=gbc_classification_2023 title="GBC 2023" title_fontsize=19 
	font="Helvetica" fontsize=17 bgcolor=white border_color=white
d.legend raster=shaded_relief title="Relief, m" title_fontsize=19 
	font="Helvetica" fontsize=17 bgcolor=white border_color=white -f
d.out.file output=GBC_2023 format=jpg --overwrite
\end{lstlisting}

\begin{lstlisting}[language=bash, caption=Support Vector Machine (SVM) Classifier script for image classification,basicstyle=\ttfamily\scriptsize,keywordstyle=\color{deepmagenta},emphstyle=\color{deepred},
stringstyle=\color{deepblue},breaklines=true,numbers=left,label={lst06}]
# training a model using r.learn.train
r.learn.train group=L_2023 training_map=training_pixels \
    model_name=SVC n_estimators=500 save_model=svc_model.gz --overwrite
# perform prediction using r.learn.predict
r.learn.predict group=L_2023 load_model=svc_model.gz 
	output=svc_classification_2014 --overwrite
# check raster categories applied to the classification output
r.category svc_classification_2023
# display the results
r.colors svc_classification_2023 color=roygbiv -e
d.mon wx2
d.rast shaded_relief
d.vect isolines color='100:93:134' width=0
d.rast svc_classification_2023
d.grid -g size=00:30:00 color=white width=0.1 fontsize=16 text_color=white
d.legend raster=svc_classification_2023 title="SVM 2023" 
	title_fontsize=19 font="Helvetica" fontsize=17 
	bgcolor=white border_color=white
d.legend raster=shaded_relief title="Relief, m" title_fontsize=19 
	font="Helvetica" fontsize=17 bgcolor=white border_color=white -f
d.out.file output=svc_2023 format=jpg --overwrite
\end{lstlisting}

{\bf Acknowledgements.} The author thanks the reviewers for reading, commenting and careful review of this paper which improved the initial version of the manuscript.

%\newsection {REFERENCES}
%\bibliographystyle{apalike}
%\bibliographystyle{plain}
%\bibliographystyle{plainnat}
\bibliography{Mauritania.bib}
%\nocite{*}

% -----------------------------------------------------------------------------------------------
% end of the AUTHOR section
% -----------------------------------------------------------------------------------------------

\end{document}